\problemname{Husbygge}
A number of eccentrics from central New York have decided that they have had enough of modern society, and want to move away.
Together they have bought a rectangular piece of land far away, and will now settle there.

The land consists of $N \times M$ cells, and at most one house can be built on a given cell.
Each cell has a value $a_{x,y}$ that describes how pleasant it is, on a scale between $0$ and $100$.

The eccentrics' goal is to get as far away as possible from everyone else, including each other.
The happiness an eccentric experiences from building their house on cell $(x,y)$ is thus $a_{x,y} \cdot d$,
where $d$ is the minimum distance to any other person.
Out of habit, the eccentrics use
\href{https://en.wikipedia.org/wiki/Taxicab_geometry}{Manhattan distance}
to measure this;
$d$ is thus defined as $\min |x - x_2| + |y - y_2|$ over all other persons' cells $(x_2, y_2)$.

The eccentrics now want your help with placing their houses optimally, so that the sum of
the happiness they experience is as high as possible. Can you help them?

\section*{Input}
The input consists of $10$ test cases, described below.

The first line contains the number $T$ ($0 \le T \le 10$), which describes the number of the test case ($0$ for sample).
The second line contains the numbers $N$, $M$ and $K$ ($1 \le N, M \le 1000$, $2 \le K \le N \cdot M$) -- the height and width of the land grid, and the number of people.
The following $N$ lines contain $M$ integers each -- the values $a_{x,y}$ ($0 \le a_{x,y} \le 100$).

\section*{Output}
Print $K$ lines with the houses' positions.
Each line should contain two numbers: first the row for the house (between $1$ and $N$), then the column (between $1$ and $M$).
Two houses may not be placed at the same position.

\section*{Scoring}
Your program will be tested on $10$ test cases, and scored based on how it compares to \emph{all other participants}.
If your submission on a given test case gets a sum of happiness of $X$, and the best sum anyone has achieved on that test case is $Y$,
then your score on that test case is $10 \cdot (X / Y)^2$, rounded to two decimal places.

Your score for a submission is the sum of the scores for the test cases, and your total score is the score for the best submission.
It is thus not possible to combine test cases between different submissions; you must submit a single program that
maximizes your score on all test cases simultaneously.

During the competition, your score will be set relative to a judge solution;
after the competition ends, all solutions will be re-scored against the best participant solutions.
If your solution is better than what the judges put together, it is therefore possible that you temporarily get more than $10$ points on a test case.
Your (temporary) total score will however be capped at $100$.

The sample test case will always give $0$ points.

\section*{Explanation of Sample}
In the sample, we want to place two houses on a $2 \times 3$ grid.
The sample solution places one of the houses in the bottom-left corner and one in the top-right corner.
Both houses' minimum distance to any other house will then be $2 + 1 = 3$,
and the sum of happiness is thus $3 \cdot 30 + 3 \cdot 50 = 240$.

If that test case had been a real test case and another participant had placed their houses
in the top-left and bottom-right corners (which would have achieved the higher happiness $270$), then
the test case would have been given $10 \cdot (240 / 270)^2 \approx 7.90$ points.

\section*{Test cases}
To avoid making the problem a guessing game, here are descriptions of the $10$ test cases:
\noindent
\begin{tabular}{| l | l | l |}
\hline
Test case & Constraints                       & Grid \\ \hline
$1$       & $N = M = 100$, $K = 1000$     & All values $a_{i,j}$ are the same. \\ \hline
$2$       & $N = M = 100$, $K = 500$      & $a_{i,j}$ is random between $0$ and $100$. \\ \hline
$3$       & $N = 200$, $M = 1$, $K = 30$  & $a_{i,j}$ is random between $0$ and $100$. \\ \hline
$4$       & $N = M = 1000$, $K = 40\,000$ & $a_{i,j}$ is random between $0$ and $100$. \\ \hline
$5$       & $N = M = 100$, $K = 20$       & $a_{i,j} = \min(100, \max(0, i + r))$, where $r$ is random between $-5$ and $5$, and $i$ is the row number, zero-indexed. \\ \hline
$6$       & $N = M = 1000$, $K = 10\,000$ & $a_{i,j} = \min(100, \max(0, \lfloor 0.101 \cdot i \rfloor + r))$, where $r$ is random between $-5$ and $5$. \\ \hline
$7$       & $N = M = 100$, $K = 500$      & $a_{i,j} = \lfloor 100 / r \rceil$, where $r$ is a random float between $1$ and $200$. \\ \hline
$8$       & $N = M = 100$, $K = 500$      & $a_{i,j} = \lfloor 100 / r^2 \rceil$, where $r$ is a random float between $1$ and $200$. \\ \hline
$9$       & $N = M = 1000$, $K = 40\,000$ & $a_{i,j} = \lfloor 100 / r^2 \rceil$, where $r$ is a random float between $1$ and $200$. \\ \hline
$10$      & $N = M = 100$, $K = 9$        & The grid consists of only $1$s with $50$ blocks of $0$s randomly placed (overlapping; most of size $10 \times 10$, some larger). \\ \hline
\end{tabular}

By random we mean \href{https://en.wikipedia.org/wiki/Continuous_uniform_distribution}{uniformly random}. $\lfloor x \rfloor$ denotes $x$ rounded down, and $\lfloor x \rceil$ denotes $x$ rounded to the nearest integer.
